\documentclass[11pt]{article}
\usepackage[utf8]{inputenc}
\usepackage[T1]{fontenc}
\usepackage{amsmath,amssymb,amsthm}
\usepackage{geometry}
\geometry{a4paper,margin=2.6cm}
\usepackage{hyperref}
\hypersetup{colorlinks=true,linkcolor=blue!50!black,urlcolor=blue!50!black}
\usepackage{booktabs}
\usepackage{xcolor}

\theoremstyle{plain}
\newtheorem{theorem}{Theorem}
\newtheorem{proposition}[theorem]{Proposition}
\theoremstyle{definition}
\newtheorem{lead}{Lead}

% verdict badges
\newcommand{\live}{\textcolor{green!45!black}{\textbf{LIVE}}}
\newcommand{\dead}{\textcolor{red!70!black}{\textbf{DEAD}}}
\newcommand{\open}{\textcolor{orange!85!black}{\textbf{OPEN}}}
\newcommand{\tool}{PTCTLMC}
\newcommand{\uppaal}{\textsc{Uppaal}}
\newcommand{\verifyta}{\texttt{verifyta}}


\title{Leads\\[2pt]\large Where the states might still be, and where they are not}
\author{PTCTLMC}
\date{August 2026}

\begin{document}
\maketitle

\noindent
This is a working document, not a paper. It is written to be read quickly and argued with.
Each section is one \emph{lead} --- an idea for exploring or storing fewer states than \uppaal{}
--- and each says plainly whether it is alive, dead, or unresolved, and why. Formal statements
appear only where a claim needs to be exact; everything else is prose.

\medskip\noindent
\textbf{The situation in one paragraph.} On five of the benchmark families (\texttt{fischer},
\texttt{arbiter}, \texttt{viking}, \texttt{twolock}, \texttt{csma}) the engine stores exactly one
zone per reachable discrete state and matches \verifyta{}'s count to the digit --- on states
explored as well as states stored. Every antichain on those families has exactly one member, so
every technique that compares two zones is provably doing nothing there. The interesting question
is therefore not ``which zone abstraction is coarser'' but ``why are we looking at that many
\emph{discrete} states at all''.

\medskip\noindent
\textbf{Summary of the leads.}

\begin{center}
\begin{tabular}{@{}llp{7.4cm}@{}}
\toprule
\S & verdict & one line \\
\midrule
1 & \live{}  & Verify a small subsystem against a bounded fake environment. \\
2 & \open{}  & A 2006 paper reports fewer states than our supposed floor. \\
3 & \dead{}  & Subsumption across different discrete states: collapse factor exactly 1. \\
4 & \dead{}  & Sweep-line: nothing to delete, and we can prove it. \\
5 & \live{}  & 41\% of expansions are wasted, and the cause is syntactic. \\
6 & \open{}  & Symmetry: the numbers are huge, the objection is fair. \\
7 & \open{}  & RWTBS: can a new bisimulation quotient more than \uppaal{}? \\
8 & \live{}  & Race a risky one-sided abstraction against the exact search. \\
\bottomrule
\end{tabular}
\end{center}

\section{Verify a small piece against a fake environment \hfill \live{}}

\textbf{The idea in one sentence.} The property \texttt{A[] not (P1.cs \&\& P2.cs)} mentions two
processes. We currently check it on all ten. Replace the other eight by a single small automaton
that is \emph{at least as permissive} as they are, and check the property on that instead.

If the replacement is genuinely at least as permissive, then anything bad that the real system can
do, the fake one can also do. So if the small system is safe, the big one is safe. The reverse is
not guaranteed --- the fake environment may invent bad behaviour that the real one cannot produce
--- which is why this is a one-way test: a ``safe'' answer is trustworthy, an ``unsafe'' answer
means we learned nothing and must fall back.

\begin{theorem}[Environment abstraction]
Let $S = C \parallel E$, where $C$ contains every component named in the safety property $\varphi$
together with anything sharing state with them, and $E$ is the rest. Let $\hat{E}$ be any timed
automaton with $E \sqsubseteq \hat{E}$ in the ready-simulation sense on the interface alphabet.
Then
\[
  C \parallel \hat{E} \models \varphi \;\Longrightarrow\; C \parallel E \models \varphi .
\]
\end{theorem}

The proof is short because someone else did the hard part: $\sqsubseteq$ is a \emph{precongruence}
for \uppaal{}'s composition operator, including shared variables and committed locations
(Berendsen \& Vaandrager, FORMATS 2008). Precongruence means the relation survives composition, so
$C \parallel E \sqsubseteq C \parallel \hat E$ follows immediately, and $\varphi$ only reads $C$,
which both sides have identically.

\paragraph{What it would buy.} On \texttt{fischer} at $N=10$, taking $\hat E$ to be one generic
replica plus the shared variable abstracted to $\{0, 1, 2, \text{other}\}$ gives roughly $256$
discrete states with $3$ clocks, against $260\,998$ with $10$ clocks. That is a factor of about
$1000$, and --- more importantly --- the abstract system does not grow with $N$ at all.

\paragraph{Is this new? No, and that is the interesting part.} The technique is quotienting, from
Andersen (LICS 1995) and Laroussinie \& Larsen (CONCUR 1995), and it was demonstrated on
\emph{exactly this model} thirty years ago: Kristoffersen, Laroussinie, Larsen, Pettersson and Yi,
TAPSOFT 1997, verified Fischer at $N = 50$ in $172$ seconds using the tool CMC, at a time when
\uppaal{} gave up around $N = 6$. Separately, Abdulla et al.\ (Nordic J.\ Computing, 2007) settled
Fischer for \emph{all} $N$ at once in $2.12$ seconds.

\paragraph{So why is nobody doing it?} Four reasons, all from the sources:

\begin{enumerate}
\item \textbf{CMC no longer exists}, and no successor tool was found.
\item \textbf{The decomposition order was chosen by hand.} The 1997 paper says so itself: the
      optimal order in which to factor out components ``needs a better understanding''. Automating
      that choice is the open problem.
\item \textbf{CMC could not handle \uppaal{} integer variables.} It modelled the shared variable as
      an explicit $(n{+}1)$-location automaton. Real models use bounded integers.
\item \textbf{The parameterized route has a sharp decidability cliff.} Reachability is decidable
      for unboundedly many identical \emph{one-clock} processes; with a controller and several
      clocks per process it is undecidable, and no cutoff theorem covers \uppaal{}-style shared
      bounded integers or binary channels.
\end{enumerate}

\paragraph{Honest risk.} The theory is proven and old. The gap is \emph{automation} --- constructing
$\hat E$ without a human. That is exactly the part nobody has done, so this lead is not ``implement
a known result'', it is ``attempt the part that stalled''. The consolation is that failure is cheap:
one abstract run that is independent of $N$, then fall back to the normal search.

\section{Is our ``floor'' actually a floor? \hfill \open{}}

Everything else in this document rests on one claim, so it deserves scrutiny.

\paragraph{The claim.} On five families we store exactly one zone per reachable discrete state.
Since you cannot store less than one thing per state you visit, and since the set of reachable
discrete states is fixed by the model, the number of stored states is a floor. \verifyta{} sits on
the same floor. Hence no zone-level cleverness can win there --- and we measured that directly:
every antichain has exactly one member, so techniques that compare pairs of zones have nothing to
compare.

\paragraph{The problem.} Ben Salah, Bozga and Maler, \emph{On Interleaving in Timed Automata},
CONCUR 2006, report Fischer at $N = 5$ in \textbf{527} symbolic states. We report $727$.
\verifyta{} reports $727$. And $727$ is the number of reachable discrete states.

So either their Fischer is a different model, or symbolic states can be fewer than discrete states
--- which would mean the floor is not a floor.

\paragraph{How that could be possible.} Their Theorem~1 says the union of all zones reached by
different interleavings of the same set of transitions is convex. That is a stronger, cleaner
statement of the thing our convex-union merge does opportunistically. If you aggregate along those
lines you are no longer storing ``one zone per discrete state'' --- you are storing one object per
\emph{equivalence class of interleavings}, and a class can span several discrete states.

\paragraph{Why this matters more than the number.} If aggregation can index fewer objects than
discrete states, then the whole framing of this document is wrong, and the right question is not
``how do we avoid visiting discrete states'' but ``what is the right unit of storage''. That would
also reframe the convex-union merge, currently our most distinctive feature, as a weak special
case of a published theorem.

\paragraph{What to do.} Reconstruct their Fischer model and count. Half a day. Two outcomes:
\begin{itemize}
\item their model differs (most likely) --- the floor claim survives and we cite the theorem as
      the general form of our merge;
\item their model is ours --- the floor claim is wrong, and this becomes the most important lead in
      the document.
\end{itemize}

Until this is settled, no paper should assert the floor.

\section{Subsumption across different discrete states \hfill \dead{}}

\textbf{The idea.} Today a stored zone can only ever kill another zone sitting at the \emph{same}
discrete state --- same locations, same integers. That is the weakest possible matching rule. If
the property only reads two processes, surely the other eight processes' locations should not have
to agree exactly for one state to subsume another?

\begin{theorem}[Componentwise-simulation subsumption]
Let $\preceq$ be a preorder on discrete states such that $\mathrm{dk}_1 \preceq \mathrm{dk}_2$
implies that every transition enabled at $\mathrm{dk}_1$ is enabled at $\mathrm{dk}_2$ with a
$\preceq$-larger target, and that $\mathrm{dk}_1, \mathrm{dk}_2$ agree on every proposition of
$\varphi$. Then $(\mathrm{dk}_1, Z)$ may be discarded whenever some stored $(\mathrm{dk}_2, Z')$
has $\mathrm{dk}_1 \preceq \mathrm{dk}_2$ and $Z \subseteq Z'$.
\end{theorem}

The relation is sound --- it is a simulation, and simulation preserves reachability of anything the
property can see.

\paragraph{And it collapses nothing.} Computed on all six families: the collapse factor is
\textbf{exactly 1}. Not ``small''. One. No pair of distinct reachable discrete states is related.

\paragraph{Why, and why this is not a benchmark artefact.} In these protocols the discrete state is
not decoration, it is the mechanism. A process sitting in \texttt{wait} rather than \texttt{idle}
changes which transitions are enabled for \emph{everyone}, because they contend for the same lock,
variable or channel. For one state to simulate another, the simulating one must enable at least as
much; here, enabling more of one process's moves systematically disables another's. The preorder
degenerates to equality.

That reasoning does not depend on the choice of benchmarks. It depends on the models being
\emph{contention protocols}, which is what timed mutual-exclusion models are. A model where
processes genuinely do not interact would collapse --- but such a model has a small product anyway,
so there is nothing to win.

\paragraph{Verdict.} Dead, and worth recording precisely because it is the obvious idea. Anyone
looking at this engine will propose it within five minutes. The answer is one measured number.

\section{Deleting states behind you \hfill \dead{}}

\textbf{The idea.} The sweep-line method (Christensen, Kristensen, Mailund) keeps a
\emph{progress measure} --- some quantity that only ever increases along transitions --- and
throws away stored states whose measure has been overtaken, because nothing can come back to them.
Peak memory falls dramatically, and memory is this engine's real wall.

\begin{theorem}[Sweep-line vacuity for cyclic reactive models]
If the reachable discrete-state graph is strongly connected, every monotone progress measure is
constant, and the sweep-line method deletes nothing.
\end{theorem}

The proof is one line: in a strongly connected graph every state reaches every other, so a measure
that never decreases along an edge cannot decrease around a cycle, hence must be equal everywhere.

\paragraph{The measurement.} Fischer's reachable discrete-state graph is a \emph{single} strongly
connected component covering $100\%$ of its states. So is every other family here. These are
reactive protocols: they return to idle and start again. There is no progress.

\paragraph{Extra reason not to pursue it.} \uppaal{} has shipped the generalised sweep-line method
since 2006. Even if it fired, it would not be a distinguishing feature.

\paragraph{Verdict.} Dead, with a proof rather than a measurement --- which is better, because it
tells you the condition under which it \emph{would} work: models with a real notion of progress,
such as pipelines or protocols with a terminating phase. None of ours has one.

\section{The 41\% of work that is thrown away \hfill \live{}}

\textbf{The observation.} We store the minimum possible number of states. We do \emph{not} expand
the minimum possible number. The difference between states explored and states stored is exactly
the number of stored zones that were later retired --- states we did the full work of expanding,
only for a better zone to arrive afterwards and make them redundant.

\begin{center}
\begin{tabular}{@{}lrr@{}}
\toprule
model & wasted expansions & share \\
\midrule
\texttt{fischer10} & $186\,600$ & $41.7\%$ \\
\texttt{twolock10} & --- & $\approx 41\%$ \\
\texttt{phil12}    & $224\,030$ & $27.9\%$ \\
\texttt{arbiter14}, \texttt{viking12}, \texttt{csma} & $0$ & $0\%$ \\
\bottomrule
\end{tabular}
\end{center}

\paragraph{Why this one is different from everything else here.} It is not benchmark-dependent, and
that is not a hope, it is measured. The delay between a state being expanded and being retired is a
\emph{constant of the model}: exactly $4$ on \texttt{fischer} at both $8$ and $10$ processes,
exactly $6$ on \texttt{twolock}, multiples of $4$ on \texttt{phil}. And the constant is
\emph{syntactic} --- it equals the length of the shortest cycle in a single process template that
returns to the same discrete key. A clock-reset side condition then predicts whether the tax exists
at all, and it predicted presence or absence correctly on $6$ of $6$ families \emph{before} running
the search.

So the precondition is checkable on any input model by looking at the templates. That is the
generality bar the project set, and this is the only lead that clears it by construction.

\paragraph{The catch, and it is a real one.} Simply holding every state back for $g$ steps is
provably useless --- the arithmetic works out so that a uniform delay costs at least as much as it
saves. This is not a small detail: it is exactly the effect Behrmann et al.\ measured in 2000 when
they distributed \uppaal{} and watched the state count \emph{rise} from $45\,001$ on one node to
$110\,634$ on seventeen. Any deferral scheme has to be selective, and the selection has to be
driven by the syntactic constant rather than by a tuned parameter.

\paragraph{The one thing that could kill it.} If the zone that retires a state is a
\emph{descendant} of that state, then the state had to be expanded to produce its own killer, and
no amount of waiting helps. A probe answered this: of $40$ sampled cases, \textbf{zero} required
the state's own expansion, and an oracle-order test attributes $86$--$89\%$ of the waste to
ordering. So the idea survives its own worst case.

\paragraph{Where parallelism comes in.} Deferral is a bet: pay latency now to avoid a subtree
later. A sequential search cannot afford the bet, because delaying its only state stalls it
completely. With many workers there is always other work, so the bet is free. This is the one place
in the whole document where extra cores buy something a sequential tool cannot have at any price.

\section{Symmetry, and why the objection to it is fair \hfill \open{}}

\textbf{The numbers first, because they are large.} The benchmark models are $N$ copies of one
template. Permuting the copies maps the system onto itself, so only one representative per orbit
needs exploring. For \texttt{fischer} with the query fixed, the number of orbits is exactly
\[
  12N^2 - 13N - 4 ,
\]
against a reachable state space that grows exponentially. That is $660$ orbits at $N=8$
($38\times$ fewer states), $1\,066$ at $N=10$ ($245\times$), and about $1\,568$ at $N=12$
($\approx 1750\times$). Exponential space, quadratic quotient.

\paragraph{The premise I used to dismiss it was wrong.} I rejected symmetry on the grounds that
``\uppaal{} already has it''. \uppaal{} has it \emph{for models written with \texttt{scalar} types}.
Our generators emit flattened templates with no scalar declarations, so \verifyta{} applies exactly
zero symmetry reduction to these inputs. Comparing our auto-detected symmetry against their
un-annotated model is therefore not a like-for-like comparison in our favour --- it is a comparison
of two different modelling styles.

\paragraph{And the objection to it is still fair.} The concern raised was that these benchmarks are
\emph{highly symmetric by construction}. Every one of them is ``$N$ identical replicas plus a shared
resource''. A technique whose payoff is proportional to how replicated the model is will look
spectacular here and do nothing on an industrial model where the components differ. That is a
benchmark-shape argument, not a soundness argument, and it is correct.

\paragraph{So what would make it honest.} Two things, and they are cheap:
\begin{enumerate}
\item Report symmetry results \emph{only} against a \uppaal{} model that also declares
      \texttt{scalar} types, so both tools reduce. Anything else is comparing annotations.
\item Include at least one asymmetric model, so the reader can see what the technique does when
      the replication is not there.
\end{enumerate}

\paragraph{Verdict.} Open, but the burden is on presentation rather than on the technique. The
$245\times$ is real; it is just measured on the shape most flattering to it. My earlier
justification for dismissing it was factually wrong, and the honest reason to be cautious is the
one that was raised: the benchmarks are too symmetric to prove anything general.

\section{RTWBS: can a new bisimulation quotient more than \uppaal{}? \hfill \open{}}

RTWBS (Relaxed Weak Timed Bisimulation) is an \emph{asymmetric} refinement relation: it insists
that a refined component preserve the timing of everything it \emph{sends}, while allowing it to be
slower about when it \emph{receives}. It is built on weak timed bisimulation, so internal
($\tau$) steps are abstracted away, and it is decided by partition refinement over the zone graphs
of the two components.

\paragraph{First, the mismatch to be honest about.} RTWBS as published answers ``does this update
refine the deployed system?'' --- a question about \emph{two} systems. Reducing the state space of
\emph{one} system is a different use. To get a quotient you take the kernel of the preorder (relate
$s$ and $t$ when each refines the other) and factor by it. That is legitimate, but it is a
different theorem from the ones in the paper, and it needs stating.

\paragraph{Why it could nevertheless be the best lead in this document.} Three reasons.

\begin{enumerate}
\item \textbf{\uppaal{} does no bisimulation quotienting at all.} It reduces by zone abstraction and
      inclusion. So any working quotient is strictly extra, not a re-implementation of something
      they have.

\item \textbf{Weak is the point.} A \emph{strong} bisimulation quotient on these models would
      collapse almost nothing, for the same reason \S3 fails: every discrete move changes what
      others can do. A \emph{weak} one abstracts internal steps, and here the property names only
      two processes --- so every move of the other $N-2$ is a candidate $\tau$. That is precisely
      the structure the relation is built to exploit, and it attacks the discrete dimension, which
      is the only dimension with slack.

\item \textbf{It composes with \S1.} The environment-abstraction lead needs a small $\hat E$ that is
      at least as permissive as the real environment. A weak-bisimulation quotient of the
      environment is a principled way to \emph{construct} one, instead of writing it by hand ---
      and hand-construction is exactly the step that stalled that line of work in 1997.
\end{enumerate}

\paragraph{The obstacle, stated plainly.} Computing a quotient of the product needs the product,
which is the thing we are trying not to build. The technique only pays if the quotient is taken
\emph{per component, before composition} --- compositional minimisation. Whether the relation is a
congruence for \uppaal{}'s composition (shared bounded integers, binary and broadcast channels,
committed locations) is the question that decides the whole lead, and it is not answered in the
paper, which composes contracts rather than quotients components.

\paragraph{Two smaller things worth noting.} The paper observes that pair validation is
``naturally parallelisable'' with an effective speedup up to the core count, which fits this
engine's setting exactly. And the relation preserves CTL$^*$ over \emph{observable events}, whereas
our properties are predicates over \emph{locations}; those need to be reconciled before any
soundness claim, most likely by making the property's locations observable and everything else
internal.

\paragraph{What would settle it, cheaply.} Take \texttt{fischer}, designate every action of
processes $3..N$ as internal, compute the weak-bisimulation quotient of a \emph{single} process
template, and count the resulting classes. If one template collapses, compositional minimisation is
worth building. If it does not, the lead dies for the same reason \S3 died, and we will know in an
afternoon.

\section{Race the risky abstraction against the exact search \hfill \live{}}

\textbf{The claim, and it is not the usual one.} Parallelism here does not make the search faster.
It changes \emph{which algorithms are usable at all}.

Every large technique in this document is sound in one direction only. Environment abstraction
(\S1): if the abstract system is safe, the real one is safe; if the abstract system is unsafe, we
learned nothing and must start over. That one-sidedness is exactly why such techniques are not in
production tools --- a sequential checker has to commit, and a wrong commitment burns the run.

With cores you do not commit. Run the cheap one-sided technique \emph{and} the exact search at the
same time. If the cheap one proves safety first, kill the exact search. If it fails, you spent one
core and the exact search was never interrupted. \textbf{Parallelism converts a risky bet into a
free one.}

\paragraph{We measured an instance of this by accident.} The liveness deadlock gate offers two
strategies that compute the same verdict:

\begin{center}
\begin{tabular}{@{}lrr@{}}
\toprule
model & gate $+$ $a_{LU}$ & exact only \\
\midrule
\texttt{fischer7} & $180\,420$ & $\mathbf{30\,994}$ \\
\texttt{arbiter6} & $\mathbf{28\,591}$ & $66\,393$ \\
\bottomrule
\end{tabular}
\end{center}

$5.8\times$ one way, $2.3\times$ the other, and \emph{which one wins cannot be known in advance} ---
deciding it is as hard as the query. Racing them yields $\min(A,B)$ on every model instead of $A$
always. It is sound because both compute the same verdict, and it is implementable immediately on a
measured spread.

\paragraph{The general version.} Put environment abstraction on the fast lane. Its cost is $O(1)$ in
$N$, so the race costs essentially one idle core, and the upside is the factor of $1000$ from \S1.
This is the specific remedy for what killed the CMC line: the decomposition order was hand-chosen
and a wrong choice wasted the run. Choose wrong in parallel and it costs nothing.

\paragraph{The counterweight, stated because it bounds the whole idea.} Parallelism also actively
\emph{hurts} state count when it is used the wrong way. Behrmann et al.\ measured \textsc{Dacapo}
rising from $45\,001$ states on one node to $110\,634$ on seventeen, purely from parallelising a
subsumption search: a state covered in one order is expanded in another. We hit the same wall from
the other side --- \verifyta{} \texttt{-o\{0,1,2\}} gives identical counts, DFS is $31.6\%$ worse on
\texttt{phil}, and \texttt{fischer8} is bit-identical from $1$ to $12$ threads \emph{only because}
strict level synchronisation is enforced.

So the rule this document ends on:

\begin{center}
\emph{Use cores for redundant, independently sound work. Never use them to relax the search order.}
\end{center}

One direction is free. The other is a tax.


\end{document}
